\section{General Results and Concrete Constructions in Model Category Theory}

This section separates results applicable to general categories from their applications to the specified generating families and model structures.
For a class of morphisms $P$, write $P^{\square}$ and ${}^{\square}P$ for its right and left lifting classes, respectively.
The Smith constructions below discharge the necessary hypotheses in particular categories; they do not assert an unconditional recognition theorem for arbitrary categories.

\begin{resultblock}
\begin{theorem}[Small object argument for finitely presentable domains]
Let $C$ have object universe $u$ and Hom universe $v$.
Suppose that $C$ is locally small relative to an auxiliary universe $w$ and has the required small colimits of that size.
Let $I$ be a small family of morphisms whose domains are finitely presentable relative to $w$.
Then $I$ satisfies the small object argument relative to $w$.
In particular, no separate cardinal condition on individual generators is needed: the relative-cell colimit condition is verified at $\aleph_0$.
\end{theorem}
\Evidence{DaggerModels.SmallObjectFiniteDomains.hasSmallObjectArgument}{DaggerModels/SmallObjectFiniteDomains.lean:56}
\Origin{A general bridge proved in this formalization and applied to Mathlib's small object argument. The three universes and the smallness, colimit, and finite-presentability hypotheses remain explicit.}
\end{resultblock}

\begin{resultblock}
\begin{theorem}[Boundary-cell presentations and lifting for simplicial sets]
In any universe $u$, every monomorphism $i:X\to Y$ of ordinary simplicial sets admits a countable relative boundary-cell presentation that retains the original $X$, $Y$, and $i$.
Let $I$ be Mathlib's actual family of boundary inclusions in all degrees.
Then
\begin{align}
I\text{-}\cell
  &= \operatorname{Mono}(\sSet), \\
I^{\square} &= \operatorname{Mono}(\sSet)^{\square}
  \subseteq \operatorname{KanFib}(\sSet).
\end{align}
The final inclusion applies to maps with arbitrary source and target.
These conclusions alone do not identify boundary right lifting with weak homotopy equivalence or with trivial fibrations.
\end{theorem}
\Evidence{DaggerModels.SSetBoundaryCells.exists_relativeCellComplex; DaggerModels.SSetBoundaryCells.boundaryCells_eq_monomorphisms; DaggerModels.SSetBoundaryCells.I_rlp_eq_monomorphisms_rlp; DaggerModels.SSetBoundaryFibration.boundary_rlp_le_fibrations}{DaggerModels/SSetBoundaryCells.lean:28,44,55; DaggerModels/SSetBoundaryFibration.lean:159}
\Origin{A boundary-cell construction and its lifting consequences proved here for ordinary simplicial sets, using Mathlib's boundary inclusions and Kan-fibration predicate.}
\end{resultblock}

\begin{resultblock}
\begin{proposition}[Promotion of a local-presentability rank]
Let $C$ have object universe $w$ and Hom universe $v$, and suppose that $C$ has pullbacks.
For regular cardinals $\kappa_0,\kappa$ in universe $u$ with $\kappa_0\leq\kappa$,
\begin{align}
C\text{ is }\kappa_0\text{-locally presentable}
  \quad\Longrightarrow\quad
C\text{ is }\kappa\text{-locally presentable}.
\end{align}
\end{proposition}
\Evidence{DaggerLocallyPresentableRankPromotion.locallyPresentable_of_le}{LocallyPresentableRankPromotion.lean:9}
\Origin{A general proposition proved here using Mathlib's strong-generator criterion. Pullbacks and the regularity of both cardinals are retained as hypotheses.}
\end{resultblock}

\begin{resultblock}
\begin{theorem}[Detection of weak equivalences by a localization]
Let $C$ be a model category in which every object is cofibrant, and let $L:C\to H$ be a localization of $C$ at its weak equivalences.
For every morphism $f$,
\begin{align}
f\text{ is a weak equivalence}
  \quad\Longleftrightarrow\quad
L(f)\text{ is an isomorphism in }H.
\end{align}
The object and Hom universes of $C$ and $H$ may be chosen independently.
\end{theorem}
\Evidence{DaggerNativeAllCofibrantLocalizationSaturation.weakEquivalence_iff_localization_isIso}{NativeAllCofibrantLocalizationSaturation.lean:32}
\Origin{A general detection theorem proved here. It assumes an existing model structure and localization; it does not construct a Bousfield localization.}
\end{resultblock}

\begin{resultblock}
\begin{proposition}[Local finite presentability of simplicial and bisimplicial sets]
For every value universe $u$, the categories of ordinary simplicial sets and bisimplicial sets are both locally finitely presentable relative to $u$:
\begin{align}
\sSet\text{ and }\operatorname{BiSSet}_u
  \quad\text{are locally finitely presentable}.
\end{align}
\end{proposition}
\Evidence{DaggerNativeSSetLocallyPresentable.locallyFinitelyPresentable; DaggerNativeBisimplicialLocallyPresentable.locallyFinitelyPresentable}{NativeSSetLocallyPresentable.lean:17; NativeBisimplicialLocallyPresentable.lean:14}
\Origin{Concrete applications in this formalization of Mathlib's results on presheaf categories and strong generators. They supply ambient-category hypotheses for the small object and Smith constructions.}
\end{resultblock}

\begin{resultblock}
\begin{theorem}[Functorial factorizations for the specified generators]
In the category $\sCat^{\dag}$ with common universe $u$, let $I_{\dag}$ be the original union of boundary generators and the object generator.
Every morphism factors functorially with left leg in ${}^{\square}((I_{\dag})^{\square})$ and right leg in $(I_{\dag})^{\square}$.
Moreover, fix an object type $O$ in universe $u$.
For each of the original generating families $I_O$ and $J_O$ in $\sCat^{\dag}_O$, every morphism $f$ has functorial factorizations
\begin{align}
f &= r_I(f)\circ\ell_I(f),
  &\ell_I(f)&\in{}^{\square}(I_O^{\square}),
  &r_I(f)&\in I_O^{\square}, \\
f &= r_J(f)\circ\ell_J(f),
  &\ell_J(f)&\in{}^{\square}(J_O^{\square}),
  &r_J(f)&\in J_O^{\square}.
\end{align}
\end{theorem}
\Evidence{DaggerModels.DaggerSmallObject.generators_hasFunctorialFactorization; DaggerFixedCategorySmallObject.I_hasFunctorialFactorization; DaggerFixedCategorySmallObject.J_hasFunctorialFactorization}{DaggerModels/DaggerSmallObject.lean:51; FixedCategorySmallObject.lean:31,35}
\Origin{Concrete applications of the general small object argument after finite presentability is proved for the relevant domains. These factorizations alone do not identify a weak-equivalence class or establish a model structure.}
\end{resultblock}

\begin{resultblock}
\begin{theorem}[A model structure on fixed-object strict dagger groupoids]
Fix a type $O$ in universe $u$.
Let $\sGpd^{\dag}_O$ be the full subcategory of $\sCat^{\dag}_O$ spanned by the objects satisfying the original degreewise invertibility condition; its morphisms are dagger-preserving functors whose object map is the identity on $O$.
This category has a model structure.
For a morphism $F:C\to D$, weak equivalences and fibrations are detected on every Hom map $F_{x,y}$, for $x,y\in O$:
\begin{align}
F\in W
  &\Longleftrightarrow
    \forall x,y\in O,\;F_{x,y}\in W_{\sSet}, \\
F\in\Fib
  &\Longleftrightarrow
    \forall x,y\in O,\;F_{x,y}\in\Fib_{\sSet}.
\end{align}
Cofibrations are the left lifting class of the specified trivial fibrations in this category.
\end{theorem}
\Evidence{DaggerOriginalFixedObjectStrictDaggerModel.strictDaggerModel; DaggerOriginalFixedObjectStrictDaggerModel.model_weakEquivalence_iff; DaggerOriginalFixedObjectStrictDaggerModel.model_fibration_iff}{OriginalFixedObjectStrictDaggerModel.lean:20,24,30,35; OriginalFixedObjectDaggerCategories.lean:87,91}
\Origin{A concrete model structure established here after the required factorizations are proved. The fixed-object category retains the original dagger enrichment and strictness. This theorem does not establish the simplicial model axioms or all of Theorem C.}
\end{resultblock}

\begin{resultblock}
\begin{theorem}[Generation of normal monomorphisms by boundaries in all degrees]
For anti-involutive simplicial sets, let $I_{\mathrm{nor}}$ be the family obtained by applying the original Part I free functor $F_{\AInv}$ to boundary inclusions in every degree.
The original class of normal monomorphisms is exactly the closure of $I_{\mathrm{nor}}$ under coproducts, pushouts, transfinite compositions, and retracts.
Indeed, every normal monomorphism admits a countable relative-cell presentation:
\begin{align}
\operatorname{NormalMono}
  = I_{\mathrm{nor}}\text{-}\cof.
\end{align}
\end{theorem}
\Evidence{DaggerModels.OriginalAntiInvolutiveNormalGeneratorRecognition.normalMonos_eq_normalGeneratorSaturation; DaggerModels.OriginalAntiInvolutiveNormalGeneratorRecognition.normalMono_mem_transfinite_pushouts}{OriginalAntiInvolutiveNormalGeneratorRecognition.lean:164,127}
\Origin{A concrete generation theorem proved here with the original normality predicate and free boundary maps in all degrees, rather than a positive-degree substitute.}
\end{resultblock}

\begin{resultblock}
\begin{theorem}[Construction of the ordinary Joyal model structure]
For ordinary $\sSet$, let $W_{\Joyal}$ be the class of weak categorical equivalences tested by the homotopy categories of internal Homs into every quasi-category $Z$.
The specified Smith data satisfy their hypotheses and give a model structure with classes
\begin{align}
\Cof_{\Joyal} &= \operatorname{Mono}(\sSet), \\
W_{\Joyal} &= \{f\mid \forall Z\text{ a quasi-category},\;
       \h(Z^f)\text{ is an equivalence of categories}\}, \\
\Fib_{\Joyal} &= (\operatorname{Mono}(\sSet)\cap W_{\Joyal})^{\square}.
\end{align}
Here $Z^f$ denotes precomposition with $f$ on the internal Hom.
The identifications of the generated left and right classes are equalities for these specified Smith data.
\end{theorem}
\Evidence{DaggerNativeOriginalJoyalSmithModelData.modelData; DaggerNativeOriginalJoyalSmithModelData.generated_left_eq; DaggerNativeOriginalJoyalSmithModelData.generated_right_eq; DaggerActualDCHRightInducedAndMainTheoremD.ordinaryJoyalModel}{NativeOriginalJoyalSmithModelData.lean:12,25,29; ActualDCHRightInducedAndMainTheoremD.lean:45}
\Origin{A concrete Smith application assembled here for ordinary simplicial sets, not a new proof of a general Smith recognition theorem. It supplies the ordinary model used in the next comparison.}
\end{resultblock}

\begin{resultblock}
\begin{theorem}[The anti-involutive Joyal model and its right-induced classes]
In the category $\sSet^{\AInv}$ of anti-involutive simplicial sets, take cofibrations to be the original normal monomorphisms and weak equivalences to be the inverse image of $W_{\Joyal}$ under the forgetful functor.
Concrete Smith data and a second functorial factorization give a model structure $\sSet^{\AInv}_{\DCH}$.
For the forgetful functor $U_{\AInv}:\sSet^{\AInv}\to\sSet$, its three classes satisfy precisely
\begin{align}
\Cof_{\DCH} &= \operatorname{NormalMono}, \\
f\in W_{\DCH}
  &\Longleftrightarrow U_{\AInv}(f)\in W_{\Joyal}, \\
f\in\Fib_{\DCH}
  &\Longleftrightarrow U_{\AInv}(f)\in\Fib_{\Joyal}.
\end{align}
These are the classes of the same concrete model used in Theorem D; its two non-Quillen conclusions are not counted again here.
Nor does this construction assert a general right-induced transfer theorem for arbitrary right adjoints.
\end{theorem}
\Evidence{DaggerActualOriginalAntiJoyalSmithModelData.modelData; DaggerActualDCHRightInducedAndMainTheoremD.model; DaggerActualDCHRightInducedAndMainTheoremD.cofibration_iff_normal; DaggerActualDCHRightInducedAndMainTheoremD.weakEquivalence_iff_underlying; DaggerActualDCHRightInducedAndMainTheoremD.fibration_iff_underlying}{ActualOriginalAntiJoyalSmithModelData.lean:20; ActualDCHRightInducedAndMainTheoremD.lean:15,22,56,64}
\Origin{A concrete Smith application in the anti-involutive category and an identification of right-induced weak equivalences and fibrations relative to the constructed ordinary Joyal model. This is the existence component of the model used in Theorem D.}
\end{resultblock}

\subsection{Foundational Results Reused from Existing Formalizations}

\begin{quote}
\raggedright
Alongside Mathlib's model-category and small-object-argument infrastructure, we reuse results ported for compatibility from the pinned version of Jo\"el Riou's
\href{https://github.com/joelriou/topcat-model-category/tree/6c0c356fea469689fe76baeb47ba773460dfacee}{\texttt{TopCatModelCategory}}.
In particular, for any Kan fibration $p:E\to B$, boundary right lifting satisfies $I^{\square}(p)\Leftrightarrow p\in W$, and ordinary $\sSet$ carries the Quillen model structure.
These are reused upstream proofs, not independent new proofs in this formalization.
\Evidence{SSet.modelCategoryQuillen.weakEquivalence_iff_of_fibration; SSet.modelCategoryQuillen.rlp_I_eq_trivialFibrations; SSet.modelCategoryQuillen.inst}{TopCatModelCategory/SSet/KeyLemma.lean:150; TopCatModelCategory/ModelCategorySSet.lean:20,51}
\end{quote}
